% ---------------------------------------------------------------
% Abstract -- one paragraph, italicised, followed by keywords.
% ---------------------------------------------------------------
\chapter*{Abstract}
\cftaddtitleline{toc}{chapter}{ABSTRACT}{\thepage}

\emph{Corrosion is a persistent threat to steel structures in the oil and
gas, marine and civil sectors, and whether a corroded component remains fit
for service depends on how deep the metal loss goes as well as how much area
it covers. Camera-based deep learning has made contactless corrosion
detection practical, but a color camera records appearance only: it can
confuse stains, coatings and shadows with corrosion, it cannot see beneath
intact paint, and it reports severity as affected area rather than loss of
material. Thermal imaging responds to defects below the surface and LiDAR
measures surface geometry, yet of the 33 corrosion-sensing studies reviewed in
this proposal, 16 rely on RGB images alone, only one combines a camera with
thermal imaging, and none fuses RGB, thermal and LiDAR data. The aim of this
study is to develop and evaluate a machine learning framework that detects
corroded regions on metal surfaces by fusing RGB, thermal and LiDAR data,
measures the corroded area and metal loss of each region from the LiDAR, and
determines what each sensor contributes to detection. A registration pipeline
will first align the three sensors, carried on an unmanned aerial vehicle,
into a common reference frame. The registered images will be stacked into a
single five-channel input to a YOLO11 detector, pretrained on general and
corrosion images, that places bounding boxes around corroded regions. An
ablation study will then compare each sensor alone, RGB--thermal, RGB--LiDAR
and the full combination, together with a blank-channel control, under
identical training conditions and a decision rule fixed in advance. A
preliminary test of this pipeline on MID-1K, a public RGB--thermal--LiDAR
person-detection dataset, ran end to end across six input configurations and
three random seeds. Early fusion reached a mean test mAP@0.5 of 0.596 against
0.572 for RGB alone, a gain too small to satisfy the decision rule, while the
same thermal images fell from 0.756 to 0.521 when scored against
RGB-aligned labels, confirming that pixel-level registration must precede
fusion. The expected contributions are a registered RGB--thermal--LiDAR
corrosion dataset and a controlled measure of which sensors are worth
deploying for corrosion detection.}

\vspace{0.8\baselineskip}

\noindent\textbf{Keywords:} \emph{Corrosion detection; multimodal sensor
fusion; early fusion; object detection; YOLO; thermal imaging; LiDAR; ablation
study.}
